%% Colors: source (DOT-MNF) and target (FCdot)
\definecolor{srccol}{RGB}{0,70,150}
\definecolor{tgtcol}{RGB}{160,40,30}
\newcommand{\src}[1]{{\color{srccol}#1}}
\newcommand{\tgt}[1]{{\color{tgtcol}#1}}
\newcommand{\run}[1]{{\color{black}#1}}
\newcommand{\mathsrc}{\everymath{\color{srccol}}\everydisplay{\color{srccol}}}
\newcommand{\mathtgt}{\everymath{\color{tgtcol}}\everydisplay{\color{tgtcol}}}
\newcommand{\mathplain}{\everymath{}\everydisplay{}}

%% Names
\newcommand{\FCdot}{\tgt{\ensuremath{\mathrm{FC}_{\mathrm{dot}}}}\xspace}
\newcommand{\DOTMNF}{\src{DOT-MNF}\xspace}
\newcommand{\Lean}{Lean~4\xspace}

%% Keywords and constructors
\newcommand{\kw}[1]{\mathbf{#1}}
\newcommand{\klet}{\kw{let}\;}
\newcommand{\kin}{\;\kw{in}\;}
\newcommand{\ev}[1]{\tgt{\mathsf{#1}}}
\newcommand{\form}[1]{\tgt{\widehat{\mathsf{#1}}}}

%% DOT types
\newcommand{\dtyp}[3]{\src{\{#1 : #2..#3\}}}
\newcommand{\dfld}[2]{\src{\{#1 : #2\}}}
\newcommand{\dmu}[2]{\src{\mu(#1.\,#2)}}
\newcommand{\dall}[3]{\src{\forall(#1 : #2)\,#3}}
\newcommand{\dsel}[2]{\src{#1.#2}}
\newcommand{\dnew}[2]{\src{\nu(#1.\,#2)}}
\newcommand{\ddeftyp}[2]{\src{\{#1 = #2\}}}
\newcommand{\ddeftrm}[2]{\src{\{#1 = #2\}}}
\newcommand{\sub}{\mathrel{\src{<:}}}

%% FCdot types
\newcommand{\sel}[2]{\tgt{#1\mathbin{\cdot}#2}}
\newcommand{\pit}[3][x]{\tgt{\Pi(#1 : #2)\,#3}}
\newcommand{\obj}[2][\self]{\tgt{\mu(#1.\,#2)}}
\newcommand{\ple}[2]{\tgt{#1 \sqsubseteq #2}}
\newcommand{\peq}[2]{\tgt{#1 \doteq #2}}
\newcommand{\phas}[1]{\tgt{\ni #1}}
\newcommand{\pbnd}[1]{\tgt{\sqsubseteq #1}}
\newcommand{\tel}{\tgt{\mathit{Tel}}}
\newcommand{\tnil}{[\,]}
\newcommand{\tcons}[2]{\tgt{#1 \triangleright #2}}
\newcommand{\tat}[3]{\tgt{#1 \ni (#2 \mapsto #3)}}
\newcommand{\weak}[1]{#1{\uparrow}}
\newcommand{\inst}[2]{#1\llbracket #2\rrbracket}
\newcommand{\self}{\tgt{\mathit{self}}}
\newcommand{\rt}[1]{\tgt{\mathit{root}(#1)}}

%% Judgments
\newcommand{\lej}[4]{\tgt{#1 \vdash #2 : #3 \le #4}}
\newcommand{\eqj}[4]{\tgt{#1 \vdash #2 : #3 \equiv #4}}
\newcommand{\hasj}[4]{\tgt{#1 \vdash #2 : #3 \ni #4}}
\newcommand{\morj}[4]{\tgt{#1 \vdash #2 : #3 \Rightarrow #4}}
\newcommand{\atomj}[3]{\tgt{#1 \vdash_{\!a} #2 : #3}}
\newcommand{\tmj}[3]{\tgt{#1 \vdash #2 : #3}}
\newcommand{\valj}[3]{\tgt{#1 \vdash_{\!v} #2 : #3}}
\newcommand{\fldj}[2]{\tgt{#1 \vdash_{\!f} #2}}
\newcommand{\storej}[2]{\tgt{\vdash #1 : #2}}
\newcommand{\contj}[4]{\tgt{#1 \vdash_{\!k} #2 : #3 \Rightarrow #4}}
\newcommand{\formj}[4]{\tgt{#1 \models #2 : #3 \le #4}}
\newcommand{\chainj}[5]{\tgt{#1 \models_{#2} #3 : #4 \le #5}}
\newcommand{\entj}[4]{\tgt{#1 \models #2 : #3 \Rightarrow #4}}
\newcommand{\entjm}[5]{\tgt{#1 \models_{#2} #3 : #4 \Rightarrow #5}}
\newcommand{\viewj}[5]{\tgt{#1 \models_{#2,#3} #4 : #5}}
\newcommand{\viewjr}[4]{\tgt{#1 \models_{#2} #3 : #4}}
\newcommand{\holej}[4]{\tgt{#1 \vdash #2 : \ple{#3}{#4}}}
\newcommand{\subj}[3]{\src{#1 \vdash #2 \sub #3}}
\newcommand{\dtyj}[3]{\src{#1 \vdash #2 : #3}}
\newcommand{\step}{\longrightarrow}
\newcommand{\steps}{\longrightarrow^{*}}
\newcommand{\erase}[1]{\lfloor #1 \rfloor}
\newcommand{\trans}[1]{{\color{black}\llbracket}\src{#1}{\color{black}\rrbracket}}
\newcommand{\resolve}[2]{\tgt{\mathit{resolve}_{#1}(#2)}}
\newcommand{\resolven}[3]{\tgt{\mathit{resolve}^{#2}_{#1}(#3)}}
\newcommand{\resolveAt}[3]{\tgt{\mathit{resolve}_{#1}^{#2}(#3)}}
\newcommand{\hnf}[4]{\tgt{#1 \vdash #2 \Downarrow_{#3} #4}}
\newcommand{\vw}[4]{\tgt{#1 \vdash #2 \Downarrow^{v}_{#3} #4}}
\newcommand{\chn}[4]{\tgt{#1 \vdash #2 \Downarrow^{c}_{#3} #4}}
\newcommand{\mst}[3]{\langle #1, #2, #3\rangle}
\newcommand{\Decl}{\src{\mathit{Decl}}}
\newcommand{\Distinct}{\src{\mathit{Distinct}}}
\newcommand{\lit}[3][\self]{\tgt{\nu(#1.\,#2;\,#3)}}
\newcommand{\ofLit}[2]{\tgt{\mathit{tel}(#1, #2)}}
\newcommand{\dom}{\mathit{dom}}

%% Rule names
\newcommand{\rn}[1]{\textsc{#1}}
\newcommand{\lean}[1]{\textup{\texttt{#1}}}

%% ---------------------------------------------------------------
%% Report 2: trees, report 1, Lean pointers
\newcommand{\treeK}{\textsf{Captures}\xspace}
\newcommand{\treeCC}{\textsf{CapturesCC}\xspace}
\newcommand{\treeCL}{\textsf{Classifiers}\xspace}
\newcommand{\Rone}{report~1~\cite{bracevac2026explicit}\xspace}
\newcommand{\leanin}[2]{\lean{#1} in #2}

%% Capture sets and the fragment
\newcommand{\cset}[1]{\{#1\}}
\newcommand{\cempty}{\{\,\}}
\newcommand{\capt}[2]{#1^{#2}}
\newcommand{\Unit}{\mathit{Unit}}
\newcommand{\CapT}[1]{\mathit{Cap}(#1)}
\newcommand{\cstar}{\ast}
\newcommand{\ann}[1]{\mathopen{\llparenthesis}#1\mathclose{\rrparenthesis}}
\newcommand{\lamc}[5]{\lambda^{#1}(#2 : #3).\,#4\;\ann{#5}}
\newcommand{\unitv}{\nu()}
\newcommand{\letc}[5]{\klet #1 = #2 \kin #3\;\ann{#4;\,#5}}
\newcommand{\frlet}[4]{\ev{let}\,#1.\,#2\;\ann{#3;\,#4}}
\newcommand{\clej}[4]{\tgt{#1 \vdash_{\!c} #2 : #3 \le #4}}
\newcommand{\ceqj}[4]{\tgt{#1 \vdash_{\!c} #2 : #3 \equiv #4}}
\newcommand{\shpj}[4]{\tgt{#1 \vdash_{\!s} #2 : #3 \le #4}}
\newcommand{\statej}[2]{\tgt{\vdash #1 : #2}}
\newcommand{\plec}[2]{\tgt{#1 \sqsubseteq_{c} #2}}
\newcommand{\peqc}[2]{\tgt{#1 \doteq_{c} #2}}
\newcommand{\uses}[1]{\mathit{uses}(#1)}
\newcommand{\insp}[1]{\mathit{insp}(#1)}
\newcommand{\annot}[1]{\mathit{annot}(#1)}
\newcommand{\roots}[2]{\mathit{roots}_{#1}(#2)}
\newcommand{\rootsn}[3]{\mathit{roots}^{#2}_{#1}(#3)}
\newcommand{\caple}[3]{#2 \preccurlyeq_{#1} #3}
\newcommand{\capeqv}[3]{#2 \approx_{#1} #3}
\newcommand{\sext}[2]{#1 \hookrightarrow #2}
\newcommand{\sextr}[3]{#2 \hookrightarrow_{#1} #3}
\newcommand{\litc}[4]{\tgt{\nu^{#1}(\self.\,#2;\,#3;\,#4)}}
\newcommand{\boxt}[1]{\tgt{\square\,#1}}
\newcommand{\unboxc}[3]{\ev{unbox}\,#1\;\ann{#2;\,#3}}

%% Source
\newcommand{\dtyuj}[4]{\src{#1;\,#2 \vdash #3 : #4}}
\newcommand{\subcj}[3]{\src{#1 \vdash #2 <:_{c} #3}}
\newcommand{\dcap}[3]{\src{\{#1 : #2..#3\}}}
\newcommand{\dbox}[1]{\src{\square\,#1}}
\newcommand{\dunbox}[2]{\src{#1 \multimap #2}}
\newcommand{\cany}{\src{\mathsf{any}}}
\newcommand{\cfresh}{\src{\mathsf{fresh}}}
\newcommand{\transa}[1]{\trans{#1}_{a}}
\newcommand{\transu}[1]{\trans{#1}_{u}}
\newcommand{\plat}{\src{P}}
\newcommand{\platctx}{\src{\Gamma_{P}}}
\newcommand{\platstore}{\src{\sigma_{P}}}
\newcommand{\tplatstore}{\llbracket\src{\sigma_{P}}\rrbracket}

%% CapturesCC
\newcommand{\topc}{\tgt{\top^{c}}}
\newcommand{\scopeb}{\mathsf{scope}}
\newcommand{\instb}[1]{\mathrel{:=}#1}
\newcommand{\lvl}[2]{\mathit{lvl}_{#1}(#2)}
\newcommand{\lvlle}[3]{#2 \trianglelefteq_{#1} #3}
\newcommand{\capsr}[3]{\mathit{caps}^{#2}_{#1}(#3)}
\newcommand{\expand}[2]{\mathit{exp}_{#1}(#2)}
\newcommand{\exc}[3]{\exists(#1 \sqsubseteq_{c} #2).\,#3}
\newcommand{\letex}[7]{\klet \langle #1, #2\rangle = #3 \kin #4\;\ann{#5;\,#6;\,#7}}
\newcommand{\mfree}[1]{\mathit{mf}(#1)}

%% Classifiers
\newcommand{\restr}[2]{#1{\restriction}#2}
\newcommand{\konly}[1]{\mathsf{only}[#1]}
\newcommand{\kexcept}[1]{\mathsf{except}[#1]}
\newcommand{\clsb}[1]{\ast_{#1}}
\newcommand{\clsof}[2]{\mathit{cls}_{#1}(#2)}
\newcommand{\kcont}[2]{#1 \mathrel{\in_{\mathsf{k}}} #2}
\newcommand{\kindj}[4]{\tgt{#1 \vdash #2 : #3 :: #4}}
\newcommand{\pkind}[2]{\tgt{#1 :: #2}}
\newcommand{\kindle}[3]{#1 \models #2 :: #3}
\newcommand{\kindsj}[3]{\src{#1 \vdash #2 :: #3}}
\newcommand{\dcapk}[2]{\src{\{#1\,{::}\,#2\}}}
